% Abstract
\vspace{1in}
\chapter*{Abstract}

IT procurement---buying, licensing, renewing and vetting hardware, software and SaaS---is still largely run through e-mail, spreadsheets and disconnected tools. The consequences are slow approvals, duplicate payments, idle licences, missed renewal deadlines and vendor risk that is assessed once and then forgotten. Techniques exist for each stage, but they are usually applied in isolation.

This report presents an IT Procurement Intelligence Platform built as four cooperating agents: Document and Vendor Intelligence, Approval and Inventory Intelligence, Contract and Risk Analysis, and a Platform agent for authentication, business rules and notifications. Each agent is an independent service that reacts to events on a Kafka-compatible bus, decides with rules or trained models, and publishes its result through a transactional outbox; long-running steps are durable Temporal workflows, and uncertain or sensitive decisions are handed to a person. Trained models are adopted only where they beat a simple baseline on held-out data. Following a comparison with five commercial suites, the platform was extended with budget control before approval, tracking of contract obligations after signature, an IT catalog, request-for-quotation events and partial goods receipts. A conversational assistant uses retrieval-augmented generation: it answers from the contracts, documents, rules and records each user is allowed to see, with citations, and answered all sixteen questions of its evaluation correctly.

On synthetic data, a TF-IDF classifier reached a macro-F1 of 1.0 against 0.912 for keyword rules; a Legal-BERT model reached F1 scores of 0.85, 0.80 and 0.91 on three clause types of the CUAD test split, against 0.50, 0.19 and 0.02 for rules; and an Isolation Forest separated anomalous from normal licence usage with a ROC-AUC of 0.954. Vendor risk is scored by a deterministic 100-point ruleset with hard overrides, in which every point traces to a check and every score can be replayed exactly; it replaced a Random Forest that had reached a ROC-AUC of 0.715 against a ceiling of 0.739. Field extraction was complete on text PDFs but reached only 49\% on image documents. On 3,000 generated invoices, the current confidence rule accepted 41\% of invoices without review, 6.8\% of them wrongly; adding purchase-order and vendor-master evidence raised the ranking quality of a calibrated gate from an AUROC of 0.74 to 0.95, but no gate could be certified at an error budget of 5\% or less, chiefly because the duplicate check misses resubmitted invoices with a new number. The implementation passes 1,282 service unit tests and end-to-end tests through the gateway. Because the data is synthetic, these figures compare methods; they do not predict performance on an organisation's own documents.

\vspace{0.5cm}
\noindent\textbf{Keywords:} IT procurement, event-driven architecture, document intelligence, human-in-the-loop, selective automation, vendor risk, anomaly detection, contract analysis, retrieval-augmented generation.
